\documentclass{book}
\usepackage{geometry}
\geometry{verbose,a4paper,tmargin=2cm,bmargin=3cm,lmargin=2cm,rmargin=3cm}
\usepackage{makeidx}
\makeindex
\usepackage{amsmath}
\usepackage{amssymb}
\usepackage{color,framed}
\usepackage{longtable}                             %for breaking tables

%%%%---- Use EPS figures and graphics with LaTeX ----
%\usepackage{graphics}
%%%%---- USE PNG or JPG figures and graphicx with PDFLaTeX ----
\usepackage{graphicx}
\usepackage[linkbordercolor={0 0.8 0.8}]{hyperref}
\hypersetup{
  pdfauthor = {Oxana Smirnova},
  pdftitle = {Extended Resource Specification Language},
  pdfsubject = {Reference manual},
  pdfkeywords = {Grid,NorduGrid,ARC,Globus,RSL},
  pdfcreator = {PDFLaTeX with hyperref package},
  pdfproducer = {PDFLaTeX}
}

\bibliographystyle{IEEEtran}

\def\efill{\hfill\nopagebreak}%
\hyphenation{Nordu-Grid}
\setlength{\parindent}{0cm}
\setlength{\FrameRule}{1pt}
\setlength{\FrameSep}{8pt}
\addtolength{\parskip}{5pt}
\renewcommand{\thefootnote}{\fnsymbol{footnote}}
\renewcommand{\arraystretch}{1.3}
\newcommand{\dothis}{\colorbox{shadecolor}}
\newcommand{\xrsl}[1]{\textbf{\sffamily #1}}
\newcommand{\versions}[1]{\vspace{-0.3cm}(#1)\par}
\newcommand{\GT}{Globus Toolkit$\mathrm{^{TM}}$}
\newcommand{\globus}{Globus\textsuperscript{\textregistered}}
\definecolor{shadecolor}{rgb}{1,1,0.6}
\definecolor{salmon}{rgb}{1,0.9,1}
\definecolor{cyan}{rgb}{0,1,1}

\begin{document}
  \def\today{\number\day/\number\month/\number\year}

  \begin{titlepage}

    \begin{tabular}{rl}
      \resizebox*{3cm}{!}{\includegraphics{ng-logo.png}}
      &\parbox[b]{2cm}{\textbf \it {\hspace*{-1.5cm}NORDUGRID\vspace*{0.5cm}}}
    \end{tabular}

    \hrulefill

    {\raggedleft NORDUGRID-MANUAL-4\par}

    {\raggedleft \today\par}

    \vspace*{2cm}

%%%%---- The title ----

    \begin{center}

    \textsc{\Large Extended Resource Specification Language}
     \Large \par \textit{Reference Manual}

    \vspace*{2cm}

%    {\centering \large Refers to ARC release series 0.8 and up \large \par}

    \end{center}

    \vspace*{2cm}

%    \abstract{
%      The paper represents an overview and a reference manual for the
%      NorduGrid Extended RSL, developed on the basis of the \globus\
%      RSL.
%    }

  \end{titlepage}
  %\thispagestyle{empty} $ $
  %\newpage
  %$\ $
  \tableofcontents
  \newpage

  \chapter{Introduction}
  \label{sec:intro}

  Computational tasks, such as data mining, data transformation, simulation of various conditions, execution of Monte Carlo algorihms and so on, are normally formalized as series of \textit{jobs} submitted to computing resources. Such computing resources can be mainframe supercomputers, high-performance computing clusters, or even pools of regular desktop PCs. In a conventional approach, applicaton experts obtain user acconts at computing centers and learn peculiarities of their \textit{batch systems} in order to formulate and submit computational jobs to the assigned resource.
  In the Grid infrastructure~\cite{grid}, experts get simultaneous access to a large number of very different computing resources, which are often widely distributed geographically. In order to abstract from the heterogeneous nature of Grid resources, a high-level job description language is necessary. The diversity of Grid resources implies special requirements for a
  proper description of a job, introducing many new options as opposed to a conventional high-performance computing center use case.

  Globus\ \index{Globus}
  Alliance~\cite{globus} developed a Grid middleware toolkit, early versions of which
  made use of the \textit{Resource Specification Language
  (RSL)}~\cite{rsl}\index{RSL} to parse job options and definitions to
  resource management systems. The NorduGrid project~\cite{nordugrid}
  \index{NorduGrid} decided to re-use RSL when they developed the Advanced Resource Connector (ARC),
  -- a
  solution~\cite{arc} for a Grid facility, suitable for
  complex tasks, like, for example, High Energy Physics
  data processing. To match the complexity of such tasks, this solution
  requires certain extensions to the RSL.

  To describe a task to be submitted to ARC-enabled resources, an
  extended version of the \globus\  RSL is used\footnote{One can also use a Grid-standard JSDL language; however, it is less versatile}. Extensions concern
  not only introduction of new attributes, but also differentiation
  between the two levels of job option specifications:

  \begin{itemize}
  \item[] \textbf{User-side RSL}, i.e., the set of attributes
    specified by a user in a job-specific file. This file is
    interpreted by a \textit{Client}~\cite{ui},
    \index{Client} and after the necessary modifications is
    passed to the ARC computing service: the \textit{Grid Manager (GM)}~\cite{gm}\index{Grid
    Manager}
  \item[] \textbf{GM-side RSL}, i.e., the set of attributes
    pre-processed by a client, and ready to be interpreted by the GM. Effectively, this is an \textit{internal} job representation of ARC.
  \end{itemize}

  A user only has to know the user-side part, and utilize it to
  describe the Grid tasks. The Grid Manager, however, uses slightly
  different notations, supplied by the client tools; therefore developers of such tools must take care of converting user-submitted RSL to the internal job description (GM-side RSL).

  In what follows, description of the NorduGrid-extended RSL,
  further denoted as \textbf{xRSL}\index{xRSL}, is given, using the
  following notations:

  \begin{tabular}{ll}
    \verb#<xxxx>#&parameter to be substituted with a corresponding string or a number\\
    \verb#[xxxx]# & optional parameter\\
    \verb#xxx|yyy|zzz# &list of possible values of a parameter\\
    \verb#-"-# &"same as above"\\
  \end{tabular}

Usage examples are given for the standard ARC command-line client (also known as \textit{User Interface})~\cite{ui}; other clients may introduce additional features, please refer to their respective documentation for details.

  \chapter{xRSL syntax and rules}
  \label{sec:remarks}

  For a complete description of \globus\  RSL, see
  reference~\cite{rsl}. xRSL uses the same syntax conventions,
  although changes the meaning and interpretation of some attributes.

  \section{Syntax}
  \label{sec:syntax}

  A Grid task is described by means of xRSL attributes, which can
  be either passed via a command-line, or, more conveniently, be
  collected in a so-called xRSL-file (suggested extension
  \textit{.xrsl}). Such a file contains a plain list of attribute-value pairs
  and boolean operators "\&" (for AND) and "\verb#|#" (for
  OR). Attribute names are case-insensitive.

  If the attributes are specified from the command line, the entire
  description must be enclosed either in single or in double quotes.
  Single quotes enclosure is completely analogous to xRSL-file usage,
  while double quotes imply standard shell expansion of the enclosed
  string. This has implications when strings contain local shell
  variables: they \textbf{will not be expanded} unless the task description is
  entered from the command line and is enclosed in double quotes.

  An attribute-value pair is a key element of specification. It
  consists of an expression that assigns one or more values to an
  attribute, and is enclosed in round brackets:

  \begin{shaded}
    \verb#(attribute="value")#\\
    \verb#(attribute="value1" "value2")# (for multi-valued attributes)
  \end{shaded}

  Certain attributes do not have assigned value; instead, they contain
  a list of values that should be put in proper correspondence with
  each other:

  \begin{shaded}
    \verb#(attribute=("value1" "value2")("value3" "value4"))#
  \end{shaded}

  In the example above, \verb#value1# and \verb#value3# are put in correspondence to
  \verb#value2# and \verb#value4# respectively, according to the
  context of the attribute.

  Values should be enclosed in quotes if they contain blank
  spaces or special characters. The \textbf{special characters} are:
  \begin{shaded}
    \verb. +   &   |   (   )   =   <   >   !   "   '   ^   #   $.%$
  \end{shaded}

  To quote a string containing special characters, you can use
  either single or double quotes. If your string, however, contains
  both such quotes, you can define any character as an own
  delimiter, by preceding it with the "carat" (\verb#^#) character:
  \verb#attribute=^*My "good" value^*# makes use of a carat-escaped
  asterisk as a delimiter.

  An xRSL job description starts with an ampersand
  ("\&") , to indicate implicit \textbf{conjunction} of all the
  attributes:
  \begin{shaded}
    \verb#&(attribute1=value1)(attribute2="value 2")...#
  \end{shaded}

  Whenever a \textbf{disjunct}-request of two or more attributes is
  needed, the following construction can be used:
  \begin{shaded}
    \verb#(|(attribute="value1")(attribute="value2")...)#
  \end{shaded}

  \begin{framed}
   Only few selected attributes (indicated further in the document) can be requested by the user multiple times, like in the disjunct request example above. Most attributes \textbf{must be unique}, i.e., appear only once in the job description document.
  \end{framed}


  In expressions, the following relational operators are allowed, in general:
  \begin{shaded}
    \verb#=   !=   >   <   >=   <=#
  \end{shaded}
  However, most attrbutes can \textbf{only be used with equality operator "="}. For few attributes (as indicated in the document), some other operators can be used as well.

  \textbf{Commented} lines should start with "(*" and be
  closed with "*)":
  \begin{shaded}
    \verb#(*attribute="value1"*)#
  \end{shaded}
  Comments can not be nested.

  \textbf{Multiple job} description in one file is realized via a standard
  \globus\  RSL multi-request operator "+", which should precede
  multiple job description:
  \begin{shaded}
    \verb#+(&(...))(&(...))(&(...))#
  \end{shaded}

  The xRSL attributes can be written in a single string, or split in
  lines arbitrary; blank spaces between and inside
  \verb#(attribute="value")# relations are ignored.

  \section{URLs}
  \label{sec:url}

  File locations in ARC can be specified both as local file
names, and as Internet standard \textit{Uniform Resource Locators
  (URL)}\index{URL}. There are also some additional URL
\textit{options} that can be used.

The following transfer protocols and metadata servers are supported:

\begin{tabular}{lp{10cm}}
   \texttt{ftp} & ordinary \textit{File Transfer Protocol (FTP)}\\
   \texttt{gsiftp} & GridFTP, the \globus\ -enhanced FTP protocol with
security, encryption, etc. developed by The Globus Alliance \cite{globus}\\
   \texttt{http} & ordinary \textit{Hyper-Text Transfer Protocol (HTTP)} with PUT and GET methods using multiple streams\\
   \texttt{https} & HTTP with SSL v3\\
   \texttt{httpg} & HTTP with \globus\  GSI\\
   \texttt{ldap} & ordinary \textit{Lightweight Data Access Protocol (LDAP)}~\cite{ldap}\\
   \texttt{rls} & \globus\  \textit{Replica Location Service (RLS)}~\cite{rls}\\
   \texttt{lfc} & LFC catalog and indexing service of EGEE gLite~\cite{glite}\\
   \texttt{srm} & Storage Resource Manager (SRM) service~\cite{srm}\\
   \texttt{file} & local to the host file name with a full path\\
   \texttt{arc} & for the Chelonia storage service, communicates with Bartenders, the path should be a Logical Name (LN)\\
\end{tabular}

An URL can be used in a standard form, i.e.
\begin{shaded}
   \verb#protocol://[host[:port]]/file#
\end{shaded}

Or, to enhance the performance, it can have additional options:
\begin{shaded}
   \verb#protocol://[host[:port]][;option[;option[...]]]/file#
\end{shaded}

For a metadata service URL, construction is the following:
\begin{shaded}
   \verb#protocol://[url[|url[...]]@]host[:port][;option[;option[...]]/#\\
   \verb#     lfn[:metadataoption[:metadataoption[...]]#
\end{shaded}

For Chelonia, the syntax is
\begin{shaded}
    \verb#arc://<LogicalName>[?BartenderURL=<URL>]#
\end{shaded}
where the BartenderURL could come from the `bartender' parameter of the client configuration file.

For the SRM service, the syntax is
\begin{shaded}
   \verb#srm://host[:port][;options]/[service_path?SFN=]file#
\end{shaded}

Versions 1.1 and 2.2 of the SRM protocol are supported. The
default \emph{service\_path} is srm/managerv2 when the server supports
v2.2, srm/managerv1 otherwise.

The URL components are:

\begin{tabular}{lp{10cm}}
   \verb#host[:port]# & Hostname or IP address [and port] of a server\\
   \verb#lfn# & Logical File Name\\
   \verb#url# & URL of the file as registered in indexing service\\
   \verb#service_path# & End-point path of the web service\\
   \verb#file# & File name with full path\\
   \verb#option# & URL option\\
   \verb#metadataoption# & Metadata option for indexing service\\
\end{tabular}

The following options are supported for location URLs:\index{URL:options}

\begin{longtable}{lp{10cm}}
   \verb#threads=<number># & specifies number of parallel
   streams to be used by GridFTP or HTTP(s,g); default value is 1,
   maximal value is 10\\
   \verb#cache=yes|no|renew|copy# & indicates whether the GM should
   cache the file; default for input files is \verb#yes#. \verb#renew#
   forces a download of the file, even if the cached copy is still valid.
   \verb#copy# forces the cached file to be copied (rather than linked) to
   the session dir, this is useful if for example the file is to be modified.\\
   \verb#readonly=yes|no# & for transfers to \verb#file://# destinations,
   specifies whether the file should be
   read-only (unmodifiable) or not; default is \verb#yes# \\
   \verb#secure=yes|no# & indicates whether the GridFTP data
   channel should be encrypted; default is \verb#no#\\
   \verb#blocksize=<number># & specifies size of
   chunks/blocks/buffers used in GridFTP or HTTP(s,g) transactions;
   default is protocol dependent\\
   \verb#checksum=cksum|md5|adler32|no# & specifies the algorithm for checksum to be
   computed (for transfer verification or provided to the indexing server). This is overridden
   by any metadata options specified (see below). If this option is
   not provided, the default for the protocol is used. \verb#checksum=no#
   disables checksum calculation.\\
   \verb#exec=yes|no# & means the file should be treated as executable\\
   \verb#preserve=yes|no# & specify if file must be uploaded to this
   destination even if job processing failed (default is \verb#no#)\\
   \verb#guid=yes|no# & make software use GUIDs instead of LFNs while
   communicating to indexing services; meaningful for \verb#rls://#
   only\\
   \verb#overwrite=yes|no# & make software try to overwrite existing
   file(s), i.e. before writing to destination, tools will try to remove
   any information/content associated with specified URL\\
   \verb#protocol=gsi|gssapi# & to distinguish between two kinds of 
   \verb#httpg#. \verb#gssapi# stands for implemention using only GSSAPI 
   functions to wrap data and \verb#gsi# uses additional headers as 
   implmented in Globus IO. The default is \verb#gssapi#. In case
   this fails, \verb#gsi# is then tried.\\
   \verb#spacetoken=<pattern># & specify the space token to be used for
   uploads to SRM storage elements supporting SRM version 2.2 or higher\\
   \verb#autodir=yes|no# & specify if before writing to specified location
   software should try to create all directories mentioned in specified
   URL. Currently this applies to FTP and GridFTP only. Default for those
   protocols is \verb#yes#\\
   \verb#tcpnodelay=yes|no# & controls the use of the TCP\_NODELAY
   socket option (which disables the Nagle algorithm). Applies to
   http(s) only. Default is \verb#no#\\ 
\end{longtable}

Local files are referred to by specifying either a location relative
to the job submission working directory, or by an absolute path (the
one that starts with "/"), preceded with a \verb#file://# prefix.

URLs also support metadata options which can be used for registering
additional metadata attributes or querying the service using metadata
attributes. These options are specified at the end of the LFN and
consist of name and value pairs separated by colons. The following
attributes are supported:

\begin{tabular}{lp{10cm}}
   \verb#guid# & GUID of the file in the metadata service \\
   \verb#checksumtype# & Type of checksum. Supported values are cksum
   (default), md5 and adler32 \\
   \verb#checksumvalue# & The checksum of the file \\
\end{tabular}

The checksum attributes may also be used to validate files that were
uploaded to remote storage.

\begin{framed}
   Examples of URLs are:\\
   \\
   \verb#http://grid.domain.org/dir/script.sh#\\
   \verb#gsiftp://grid.domain.org:2811;threads=10;secure=yes/dir/input_12378.dat#\\
   \verb#ldap://grid.domain.org:389/lc=collection1,rc=Nordugrid,dc=nordugrid,dc=org#\\
   \verb#rls://gsiftp://se.domain.org/datapath/file25.dat@grid.domain.org:61238/myfile02.dat#$^1$\\
   \verb#file:///home/auser/griddir/steer.cra#\\
   \verb#lfc://srm://srm.domain.org/griddir@lfc.domain.org/user/file1:guid=\# \\
   \verb#    bc68cdd0-bf94-41ce-ab5a-06a1512764dc:checksumtype=adler32:checksumvalue=12345678#$^2$\\
   \verb#lfc://lfc.domain.org;cache=no/:guid=bc68cdd0-bf94-41ce-ab5a-06a1512764d#$^3$\\
\end{framed}

$^1$This is a destination URL. The file will be copied to the GridFTP
server at \texttt{se.domain.org} with the path \texttt{datapath/file25.dat} and
registered in the RLS indexing service at \texttt{grid.domain.org} with the LFN
\texttt{myfile02.dat}.

$^2$This is a destination URL. The file will be copied to
\texttt{srm.domain.org} at the path \texttt{griddir/file1} and registered to the LFC
service at \texttt{lfc.domain.org} with the
LFN \texttt{/user/file1}. The checksum will be compared to what is
reported by the SRM service after the transfer. The given GUID and
checksum attributes will also be registered in the LFC.

$^3$This is a source URL. The file is registered in the LFC service at
\texttt{lfc.domain.org} with the given GUID and can be copied or queried by
this URL.


  \chapter{Attributes}
  \label{sec:attributes}

  Most of the job description attributes \index{attributes} introduced originally by \globus\ RSL~1.0 are
  supported, some with modifications as indicated in this
  document. Many new attributes are introduced by ARC, of which some are
  to be specified in the user's script, and others are internal for
  the GM (are added or modified by client tools).

  Attribute names are case-insensitive, although assigned values may
  well be case-sensitive, if they represent file names, environment
  variables etc..

  It is possible to use unsupported attributes in job description. Standard ARC client submission commands (\texttt{ngsub} and \texttt{ngresub}) must be used with a command line option ``-U'' in order to accept unknown attributes. Without this command line option, the client tool will concider job description invalid if it contains unsupported attributes.

  \section{User-side attributes}
  \label{sec:user}

  The following attributes \index{attributes:user-side} can be
  specified in a user's xRSL script. Some have to be modified by the client tool
  before being passed to the GM. If this is the case, the corresponding modified GM input is described in this document as well.

    \hspace*{0.5cm}\phantomsection
  \index{XRSL:attribute:executable}\addcontentsline{toc}{subsection}{executable}
  \begin{shaded}
    \xrsl{executable}
  \end{shaded}
  \versions{ARC 0.3, ARC 0.4, ARC 0.5, ARC 0.6}
  \begin{tabular}{lp{13cm}}
    User input:&\verb#(executable=<string>)#\\
    GM input:& for ARC $\leq$ 0.5.25:  \verb#(executable="/bin/echo")#\\
    Example:&\verb#(executable="local_to_job.exe")#\\
  \end{tabular}

  The executable to be submitted as a main task to a Local Resource
  Management System (LRMS).

  \begin{tabular}{llp{10cm}}
    \hspace*{1cm}&\texttt{string}&file name (including path), local to the computing
    element (CE)\\
  \end{tabular}

  Executable is a file that has to be executed as the main process
  of the task. It could be either a pre-compiled binary, or a
  script. Users may transfer their own executables, or use the ones
  known to be already installed on the remote system (CE).

  If an executable has to be transferred to the destination site (CE) from some source, it
  has to be specified in the \texttt{inputFiles} list. If it is not specified in \texttt{inputFiles},
  the source is expected to be local to the user (client) and
  will be added as such to the \texttt{inputFiles} list by the UI.

  If the file name starts with a leading slash ("\verb#/#"), it is considered
  to be \textbf{the full path to the executable at the destination site (CE)}; otherwise
  the location of the file is \textbf{relative} to the session
  directory (where job input and files are stored). 

  If the xRSL string is entered from the command line and is enclosed
  in double quotes, standard shell expansion of variables takes place.
  That is, if the file name
  contains an environment variable ("\verb#$...#"), %$
  the value of this variable is resolved locally, but if the name
  itself is also
  enclosed in double quotes, it will be resolved at the remote
  computing element:\\
  \verb#(executable=$ROOT_DIR/myprog.exe)# -- \verb#$ROOT_DIR# is resolved
  locally (\textit{will cause errors if the path does not exist at the
  execution machine})\\
  \verb#(executable="$ROOT_DIR/myprog.exe")# -- \verb#$ROOT_DIR# will be
  resolved remotely\\

  For ARC $\leq$ 0.5.25, UI created a dummy executable attribute for internal GM use,
  with e.g., "\verb#/bin/echo#" as an argument (unless the
  specified executable itself started with "\verb#/#"). The actual
  executable was transferred by the UI to
  \texttt{arguments}, \texttt{inputFiles} and \texttt{executables} lists
  in the pre-processed xRSL script.

  \hspace*{0.5cm}\phantomsection
  \index{XRSL:attribute:arguments}\addcontentsline{toc}{subsection}{arguments}
  \begin{shaded}
    \xrsl{arguments}
  \end{shaded}
  \versions{ARC 0.3, ARC 0.4, ARC 0.5, ARC 0.6}
  \begin{tabular}{lp{13cm}}
    User input:&\verb#(arguments=<string> [string] ... )#\\
    GM input:&\verb#(arguments=<executable> <string> [string] ... )#\\
    Example:&\verb#(arguments="10000" $(ATLAS)/input.dat)#\\ %$
  \end{tabular}

  List of the arguments for the executable.

  \begin{tabular}{llp{10cm}}
    \hspace*{1cm}&\texttt{string}&an argument\\
    \hspace*{1cm}&\texttt{executable}&the executable to be run by LRMS, taken by the
    UI from the user-specified \texttt{executable} attribute\\
  \end{tabular}

  For ARC $\leq$ 0.5.25, UI added the name specified in \texttt{executable} to the 
  \texttt{arguments} list in the pre-processed xRSL script.

  \hspace*{0.5cm}\phantomsection
  \index{XRSL:attribute:inputFiles}\addcontentsline{toc}{subsection}{inputFiles}
  \begin{shaded}
    \xrsl{inputFiles}
  \end{shaded}
  \versions{ARC 0.3, ARC 0.4, ARC 0.5, ARC 0.6}
  \begin{tabular}{lp{12cm}}
    User input:&\verb#(inputFiles=(<filename> <location>) ... )#\\
    GM input:&\verb#(inputFiles=(<filename> <URL>)#\\
    &\verb#        (<filename> [size][.checksum]) ... )#\\
    Example:&\verb#(inputFiles=("local_to_job" "gsiftp://se1.lu.se/p1/remote_one")#\\
    &\verb#        ("local_to_job.dat" "/scratch/local_to_me.dat")#\\
    &\verb#        ("same_name_as_in_my_current_dir" ""))#\\
  \end{tabular}

  List of files to be copied to the computing element before the
  execution.

  \begin{tabular}{llp{10cm}}
    \hspace*{1cm}&\texttt{filename} & file name, local to the computing element and
    always relative to the session directory\\
    \hspace*{1cm}&\texttt{location} & location of the file (\texttt{gsiftp},
    \texttt{https}, \texttt{ftp}, \texttt{http} URLs, or a path, local
    to the submission node). If void ("", use the quotes!), the input file is taken
    from the submission directory.\\
    \hspace*{1cm}&\texttt{URL} & URL of the file (see Section~\ref{sec:url})\\
    \hspace*{1cm}&\texttt{size} & file size in bytes\\
    \hspace*{1cm}&\texttt{checksum} & file checksum (as returned by \texttt{cksum})\\
  \end{tabular}

  If the list does not contain the standard input file (as specified
  by \texttt{stdin}) and/or the executable file (as specified by
  \texttt{executable} if the given name), the UI appends these files
  to the list. If the \verb#<location># is a URL, it is passed by the
  UI to the GM without changes; if it is a local path (or void, ""),
  the UI converts it to \verb#<size># and \verb#<checksum># and
  uploads those files to the execution cluster.

  \begin{framed}
  Please note that the \texttt{inputFiles} attribute is not meant to
  operate with directories: you must specify a pair
  \verb#("<local_to_job>" "[source]")# for each file.
  \end{framed}

  For ARC $\leq$ 0.5.25, UI added the name specified in \texttt{executable} to the 
  \texttt{inputFiles} list in the pre-processed xRSL script.

  \hspace*{0.5cm}\phantomsection
  \index{XRSL:attribute:executables}\addcontentsline{toc}{subsection}{executables}
  \begin{shaded}
    \xrsl{executables}
  \end{shaded}
  \versions{ARC 0.3, ARC 0.4, ARC 0.5, ARC 0.6}
  \begin{tabular}{lp{13cm}}
    User input:&\verb#(executables=<string> [string] ...)#\\
    GM input:&\verb#-"-#\\
    Example:&\verb#(executables="myscript.sh" "myjob.exe")#\\
  \end{tabular}

  List of files from the \texttt{inputFiles} set, which will be given
  executable permissions.

  \begin{tabular}{llp{10cm}}
    \hspace*{1cm}&\texttt{string}& file name, local to the computing element and
    relative to the session directory\\
  \end{tabular}

  If the executable file (as specified in \texttt{executable} and is
  relative to the session directory)  is not listed, it will be added
  to the list by the UI.

  \hspace*{0.5cm}\phantomsection
  \index{XRSL:attribute:cache}\addcontentsline{toc}{subsection}{cache}
  \begin{shaded}
    \xrsl{cache}
  \end{shaded}
  \versions{ARC $\geq$ 0.3.34, ARC 0.4, ARC 0.5, ARC 0.6}
  \begin{tabular}{lp{13cm}}
    User input:&\verb#(cache="yes"|"no")#\\
    GM input:&\verb#-"-#\\
    Example:&\verb#(cache="yes")#\\
  \end{tabular}

  Specifies whether input files specified in the \texttt{inputFiles}
  should be placed by default in the cache or not. This does not
  affect files described by the \texttt{executables}, which will be
  placed in the session directory always.

  If not specified, default value is "yes".

  \begin{framed}
    Cached files can not be modified by jobs by default. If your job has to 
    modify input files, please specify \texttt{(cache="no")} or use
    \texttt{(readonly="no")} URL option for those files.
  \end{framed}

  \hspace*{0.5cm}\phantomsection
  \index{XRSL:attribute:outputFiles}\addcontentsline{toc}{subsection}{outputFiles} 
  \begin{shaded}
    \xrsl{outputFiles}
  \end{shaded}
  \versions{ARC 0.3, ARC 0.4, ARC 0.5, ARC 0.6}
  \begin{tabular}{lp{13cm}}
    User input:&\verb#(outputFiles=(<string> <URL>) ... )#\\
    GM input:&\verb#-"-#\\
    Example:&\verb#(outputFiles=("local_to_job.dat" "gsiftp://se1.uo.no/stored.dat")#\\
    &\verb#    ("local_to_job.dat" "rc://rc1.uc.dk/set1/stored_and_indexed.dat")#\\  
    &\verb#    ("local_to_job_dir/" ""))#\\  
  \end{tabular}

  List of files to be retrieved by the user or uploaded by the GM and
  indexed (registered) in a data indexing service, e.g. \globus\  RLS or Fireman.

  \begin{tabular}{llp{10cm}}
    \hspace*{1cm}&\texttt{string}& file name, local to the
    \textit{Computing Element (CE)}\index{Computing Element}. If this
    string ends with a backslash "/" and \verb#<URL># is empty, the
    entire directory is being kept at the execution site. If however
    this string ends with a backslash "/" but the \verb#<URL>#
    starts with \texttt{gsiftp} or \texttt{ftp}, the whole directory
    is transfered to the destination.\\
    \hspace*{1cm}&\texttt{URL} & URL of the remote file (see
    Section~\ref{sec:url}); if void ("", use the quotes!), the file is
    kept for manual retrieval.  Note that this can not be a local
    \texttt{file://} URL.\\
  \end{tabular}

  Using a Replica Catalog pseudo-URL (see Section~\ref{sec:url}), you should make
  sure that the location is already defined in the Replica Catalog. If more
  than one such location is specified, only those found in the Replica Catalog
  for this collection will be used. Grid Manager will store output files in
  \textbf{one} location only: the first randomly picked location that
  exists. If no locations are specified, all those found in the
  Replica Catalog for this collection will be tried.

  If in a \verb#rc://# pseudo-URL the server component \texttt{host[:port]/DN} is not
  specified, server specified in the \texttt{replicaCollection} attribute will be used.

  If the list does not contain standard output, standard error
  file names and GM log-files directory name (as specified by \texttt{stdout},
  \texttt{stderr} and \texttt{gmlog}), the UI appends these
  items to the \texttt{outputFiles} list. If the \verb#<URL># is not specified (void,
  "", use the quotes!), files will be kept on SE and should be downloaded 
  by the user via the UI. If specified name of file ends with "/", the 
  entire directory is kept.

  A convenient way to keep the entire job directory at the remote site
  for a manual retrieval is to specify \texttt{(outputfiles=("/" ""))}.

  In some cases, the list of output files may only be known after the
  job has completed. ARC allows a user to specify a list of output
  files dynamically in a file or files in the session directory as
  part of their job. The file(s) containing the output file
  information can be defined in the xRSL script as the path to the
  file relative to the session directory preceeded by '@'. The format
  of these files is lines of 2 values separated by a space. The first
  value contains name of the output file relative to the session
  directory and the second value is a URL to which the file will be
  uploaded.

  \begin{tabular}{lp{13cm}}
    Example:&\verb#(outputFiles=("@output.files" "")#\\
    & \emph{output.files} is generated by the user and contains \\
    &\verb#   file1 gsiftp://grid.domain.org/file1#\\
    &\verb#   file2 gsiftp://grid.domain.org/file2#\\
  \end{tabular}

  After the job completes, the file output.files in the session
  directory will be read and any files described within will be
  uploaded to the given URLs.

  \begin{framed}
    Please note that the \texttt{outputFiles} attribute is not meant to
    operate with directories: you must specify a pair
    \verb#("<local_to_job>" "[destination]")# for each file. One
    exception is when you want to preserve an entire directory at the
    remote computer for later \textbf{manual} download via
    \texttt{ngget}. In that case, simply add the trailing backslash
    "/" at the end of the remote directory name. You can not upload a
    directory to a URL location, only to your local disk.
  \end{framed}

  \hspace*{0.5cm}\phantomsection
  \index{XRSL:attribute:cpuTime}\addcontentsline{toc}{subsection}{cpuTime}
  \begin{shaded}
    \xrsl{cpuTime}
  \end{shaded}
  \versions{ARC $\geq$ 0.3.17\footnote{Was called \texttt{maxCpuTime} for ARC $\leq$ 0.3.16.}, ARC 0.4, ARC 0.5, ARC 0.6}
  \begin{tabular}{lp{13cm}}
    User input:&\verb#(cpuTime=<time>)#\\
    GM input:&\verb#(cpuTime=<tttt>)#\\
    Example:&\verb#(cpuTime="240")#\\
  \end{tabular}

  Maximal CPU time request for the job. 

  \begin{tabular}{llp{10cm}}
    \hspace*{1cm}&\texttt{time} & time (minutes)\\
    \hspace*{1cm}&\texttt{tttt} & time converted by the UI from \texttt{time} to 
    minutes (ARC $\leq$ 0.5.26) or seconds (ARC $\geq$ 0.5.27)\\
  \end{tabular}
  
  If only number is specified, the time is assumed to be minutes (ARC
  $\leq$ 0.5.26) or seconds (ARC $\geq$ 0.5.27). Otherwise, a free
  format is accepted, i.e., any of the following will be interpreted
  properly (make sure to enclose such strings in quotes!):

  \begin{tabular}{ll}
    \hspace*{1cm}&\texttt{"1 week"}\\
    \hspace*{1cm}&\texttt{"3 days"}\\
    \hspace*{1cm}&\texttt{"2 days, 12 hours"}\\
    \hspace*{1cm}&\texttt{"1 hour, 30 minutes"}\\
    \hspace*{1cm}&\texttt{"36 hours"}\\
    \hspace*{1cm}&\texttt{"9 days"}\\
    \hspace*{1cm}&\texttt{"240 minutes"}\\
  \end{tabular}
 
  If both \texttt{cpuTime} and \texttt{wallTime} are specified, the UI converts them both to 
  seconds. \texttt{cpuTime} can not be specified together with \texttt{gridTime}
  or \texttt{benchmarks}.
 
  \begin{framed}
    This attribute should be used to direct the job to a system with
    sufficient CPU resources, typically, a batch queue with the
    sufficient upper time limit. Jobs exceeding this maximum most
    likely will be \textbf{terminated} by remote systems! If
    time limits are not specified, the limit is not set and jobs
   can run as long as the system settings allow (note that in this
    case you can not avoid queues with too short time limits).
  \end{framed}

  \hspace*{0.5cm}\phantomsection
  \index{XRSL:attribute:wallTime}\addcontentsline{toc}{subsection}{wallTime}
  \begin{shaded}
    \xrsl{wallTime}
  \end{shaded}
  \versions{ARC $\geq$ 0.5.27, ARC 0.6}
  \begin{tabular}{lp{13cm}}
    User input:&\verb#(wallTime=<time>)#\\
    GM input:&\verb#(wallTime=<tttt>)#\\
    Example:&\verb#(wallTime="240")#\\
  \end{tabular}

  Maximal wall clock time request for the job. 

  \begin{tabular}{llp{10cm}}
    \hspace*{1cm}&\texttt{time} & time (minutes)\\
    \hspace*{1cm}&\texttt{tttt} & time converted by the UI from \texttt{time} to seconds\\
  \end{tabular}
  
  If only number is specified, the time is assumed to be
  seconds. Otherwise, a free format is accepted, i.e., any of the
  following will be interpreted properly (make sure to enclose such
  strings in quotes!):

  \begin{tabular}{ll}
    \hspace*{1cm}&\texttt{"1 week"}\\
    \hspace*{1cm}&\texttt{"3 days"}\\
    \hspace*{1cm}&\texttt{"2 days, 12 hours"}\\
    \hspace*{1cm}&\texttt{"1 hour, 30 minutes"}\\
    \hspace*{1cm}&\texttt{"36 hours"}\\
    \hspace*{1cm}&\texttt{"9 days"}\\
    \hspace*{1cm}&\texttt{"240 minutes"}\\
  \end{tabular}
 
  If both \texttt{cpuTime} and \texttt{wallTime} are specified, the UI converts them both to 
  seconds. \texttt{wallTime} can not be specified together with \texttt{gridTime}
  or \texttt{benchmarks}. If only \texttt{wallTime} is specified, but not \texttt{cpuTime},
  the corresponding  \texttt{cpuTime} value is evaluated by the UI and added to the job description.
 
  \begin{framed}
    This attribute should be used to direct the job to a system with
    sufficient CPU resources, typically, a batch queue with the
    sufficient upper time limit. Jobs exceeding this maximum most
    likely will be \textbf{terminated} by remote systems! If 
    time limits are not specified, the limit is not set and jobs
    can run as long as the system settings allow (note that in this
    case you can not avoid queues with too short time limits).
  \end{framed}

  \hspace*{0.5cm}\phantomsection
  \index{XRSL:attribute:gridTime}\addcontentsline{toc}{subsection}{gridTime}
  \begin{shaded}
    \xrsl{gridTime}
  \end{shaded}
  \versions{ARC $\geq$ 0.3.31, ARC 0.4, ARC 0.5, ARC 0.6}
  \begin{tabular}{lp{13cm}}
    User input:&\verb#(gridTime=<time>)#\\
    GM input:& none\\
    Example:&\verb#(gridTime="2 h")#\\
  \end{tabular}

  Maximal CPU time request for the job scaled to the 2.8~GHz
  Intel\textsuperscript{\textregistered}
  Pentium\textsuperscript{\textregistered}~4 processor.

  \begin{tabular}{llp{10cm}}
    \hspace*{1cm}&\texttt{time} & time (minutes)\\
  \end{tabular}
  
  The attribute is completely analogous to \texttt{cpuTime}, except
  that it will be recalculated to the actual CPU time request for each
  queue, depending on the published processor clock speed.
  
  \texttt{gridTime} can not be specified together with \texttt{cpuTime}
  or \texttt{wallTime}. If only \texttt{gridTime} is specified, but not \texttt{cpuTime},
  the corresponding  \texttt{cpuTime} value is evaluated by the UI and added to the job description.

  \hspace*{0.5cm}\phantomsection
  \index{XRSL:attribute:benchmarks}\addcontentsline{toc}{subsection}{benchmarks}
  \begin{shaded}
    \xrsl{benchmarks}
  \end{shaded}
  \versions{ARC $\geq$ 0.5.14, ARC 0.6}
  \begin{tabular}{lp{13cm}}
    User input:&\verb#(benchmarks=(<string> <value> <time>) ... )#\\
    GM input:&\\
    Example:&\verb#(benchmarks=("mybenchmark" "10" "1 hour, 30 minutes"))#\\
  \end{tabular}

  Evaluate a job's \texttt{cpuTime} based on benchmark values.

  \begin{tabular}{llp{10cm}}
    \hspace*{1cm}&\texttt{string} & benchmark name\\
    \hspace*{1cm}&\texttt{value} & benchmark value of reference machine\\
    \hspace*{1cm}&\texttt{time} & the \texttt{cpuTime} the job requires on the reference machine\\
  \end{tabular}
  
  \texttt{benchmarks} can not be specified together with \texttt{cpuTime}
  or \texttt{wallTime}. If only \texttt{benchmarks} is specified, but not \texttt{cpuTime},
  the corresponding  \texttt{cpuTime} value is evaluated by the UI and added to the job description.

  \hspace*{0.5cm}\phantomsection
  \index{XRSL:attribute:memory}\addcontentsline{toc}{subsection}{memory}
  \begin{shaded}
    \xrsl{memory}
  \end{shaded}
  \versions{ARC $\geq$ 0.3.17\footnote{Was called \texttt{maxMemory} for ARC $\leq$ 0.3.16}, ARC 0.4, ARC 0.5, ARC 0.6}
  \begin{tabular}{lp{13cm}}
    User input:&\verb#(memory=<integer>)#\\
    GM input:&\verb#-"-#\\
    Example:&\verb#(memory>="500")#\\
  \end{tabular}

  Memory required for the job.

  \begin{tabular}{llp{10cm}}
    \hspace*{1cm}&\texttt{integer} & size (Mbytes)\\
  \end{tabular}

  \begin{framed}
    Just like \texttt{cpuTime}, this attribute should be used to
    direct a job to a resource with a sufficient capacity. Jobs
    exceeding this memory limit will most likely be \textbf{terminated}
    by the remote system.
  \end{framed}

  \hspace*{0.5cm}\phantomsection
  \index{XRSL:attribute:disk}\addcontentsline{toc}{subsection}{disk}
  \begin{shaded}
    \xrsl{disk}
  \end{shaded}
  \versions{ARC $\geq$ 0.3.17\footnote{Was called \texttt{maxDisk} for ARC $\leq$ 0.3.16}, ARC 0.4, ARC 0.5, ARC 0.6}
  \begin{tabular}{lp{13cm}}
    User input:&\verb#(disk=<integer>)#\\
    GM input:&none\\
    Example:&\verb#(disk="500")#\\
  \end{tabular}

  Disk space required for the job.

  \begin{tabular}{llp{10cm}}
    \hspace*{1cm}&\texttt{integer}  & disk space, Mbytes\\
  \end{tabular}

  \begin{framed}
    This attribute is used at the job submission time to find a system
    with sufficient disk space. However, it \textbf{does not
    guarantee} that this space will be available at the end of the
    job, as most known systems do not allow for disk space
    allocation. Eventually, a remote system can terminate a job that
    exceeds the requested disk space.
  \end{framed}

  \hspace*{0.5cm}\phantomsection
  \index{XRSL:attribute:runTimeEnvironment}\addcontentsline{toc}{subsection}{runTimeEnvironment}
  \begin{shaded}
    \xrsl{runTimeEnvironment}
  \end{shaded}
  \versions{ARC 0.3, ARC 0.4, ARC 0.5, ARC 0.6}
  \begin{tabular}{lp{13cm}}
    User input:&\verb#(runTimeEnvironment=<string>)(runTimeEnvironment=<string>)#\\
    GM input:&For ARC $\leq$ 0.5.39: \verb#(runTimeEnvironment="<string>" "<string>")#\\
    Example:&\verb#(runTimeEnvironment>="APPS/HEP/ATLAS-10.0.1")#\\
  \end{tabular}

  Required runtime environment.

  \begin{tabular}{llp{10cm}}
    \hspace*{1cm}&\texttt{string} & environment name\\
  \end{tabular}

  The site to submit the job to will be chosen by the UI among those
  advertising specified runtime environments. Before starting the job,
  the GM will set up environment variables and paths according to
  those requested. Runtime environment names are defined by Virtual 
  Organizations, and tend to be organized in name spaces.

  To request several environments, repeat the attribute string:\\
  \verb#(runTimeEnvironment="ENV1")(runTimeEnvironment="ENV2")# etc.   
  In ARC $\leq$ 0.5.39, before being submitted to the GM, a grouping was done by the UI, such that GM input was,
  e.g., \verb#(runTimeEnvironment="ENV1" "ENV2")#.

  To make a disjunct-request, use a boolean expression:\\
  \verb#(|(runTimeEnvironment="env1")(runTimeEnvironment="env2"))#.

  You can use "\verb#>=#" or "\verb#<=#" operators: job will be submitted to
  any suitable site that satisfies such requirements, and among the available at the sites
  runtime environments, the highest version satisfying a requirement will be requested
  in the pre-processed xRSL script.

  Runtime environment string interpretation is
  case-insensitive.  If a runtime environment string consists of a name
  and a version number, a partial specification is possible: it is sufficient
  to request only the name and use "\verb#>#" or "\verb#>=#" operators to select the highest version.

  \hspace*{0.5cm}\phantomsection
  \index{XRSL:attribute:middleware}\addcontentsline{toc}{subsection}{middleware}
  \begin{shaded}
    \xrsl{middleware}
  \end{shaded}
  \versions{ARC 0.3, ARC 0.4, ARC 0.5, ARC 0.6}
  \begin{tabular}{lp{13cm}}
    User input:&\verb#(middleware=<string>)#\\
    GM input:&\verb#-"-#\\
    Example:&\verb#(middleware="nordugrid-arc-0.5.99")#\\
  \end{tabular}

  Required middleware.

  \begin{tabular}{llp{10cm}}
    \hspace*{1cm}&\texttt{string} & Grid middleware name.\\
  \end{tabular}

  The site to submit the job to will be chosen by the UI among those
  advertising specified middleware. Usage is identical to that of the
  \texttt{runTimeEnvironment}. Use the "\verb#>=#" operator to
  request a version "equal or higher". Request
  \texttt{(middleware="nordugrid-arc")} defaults to
  \verb#(middleware>="nordugrid-arc-0.0.0.0")#.

  \hspace*{0.5cm}\phantomsection
  \index{XRSL:attribute:opsys}\addcontentsline{toc}{subsection}{opsys}
  \begin{shaded}
    \xrsl{opsys}
  \end{shaded}
  \versions{ARC $\geq$ 0.3.21, ARC 0.4, ARC 0.5, ARC 0.6}
  \begin{tabular}{lp{13cm}}
    User input:&\verb#(opsys=<string>)#\\
    GM input:&\verb#-"-#\\
    Example:&\verb#(opsys="FC3")#\\
  \end{tabular}

  Required operating system.

  \begin{tabular}{llp{10cm}}
    \hspace*{1cm}&\texttt{string} & Operating system name and version.\\
  \end{tabular}

  The site to submit the job to will be chosen by the UI among those
  advertising specified operating system. Usage is identical to that of
  \texttt{runTimeEnvironment} and \texttt{middleware}. Use the "\verb#>=#" operator to
  request a version "equal or higher".

  \hspace*{0.5cm}\phantomsection
  \index{XRSL:attribute:stdin}\addcontentsline{toc}{subsection}{stdin}
  \begin{shaded}
    \xrsl{stdin}
  \end{shaded}
  \versions{ARC 0.3, ARC 0.4, ARC 0.5, ARC 0.6}
  \begin{tabular}{lp{13cm}}
    User input:&\verb#(stdin=<string>)#\\
    GM input:&\verb#-"-#\\
    Example:&\verb#(stdin="myinput.dat")#\\
  \end{tabular}

  The standard input file.

  \begin{tabular}{llp{10cm}}
    \hspace*{1cm}&\texttt{string} & file name, local to the computing element\\
  \end{tabular}

  The standard input file should be listed in the \texttt{inputFiles}
  attribute; otherwise it will be forced to that list by the UI.

  \hspace*{0.5cm}\phantomsection
  \index{XRSL:attribute:stdout}\addcontentsline{toc}{subsection}{stdout}
  \begin{shaded}
    \xrsl{stdout}
  \end{shaded}
  \versions{ARC 0.3, ARC 0.4, ARC 0.5, ARC 0.6}
  \begin{tabular}{lp{13cm}}
    User input:&\verb#(stdout=<string>)#\\
    GM input:&\verb#-"-#\\
    Example:&\verb#(stdout="myoutput.txt")#\\
  \end{tabular}

  The standard output file.

  \begin{tabular}{llp{10cm}}
    \hspace*{1cm}&\texttt{string} & file name, local to the computing element and
    relative to the session directory.\\
  \end{tabular}

  The standard output file should be listed in the
  \texttt{outputFiles} attribute; otherwise it will be forced to that
  list by the UI. If the standard output is not defined, UI assigns a
  name.

  \hspace*{0.5cm}\phantomsection
  \index{XRSL:attribute:stderr}\addcontentsline{toc}{subsection}{stderr}
  \begin{shaded}
    \xrsl{stderr}
  \end{shaded}
  \versions{ARC 0.3, ARC 0.4, ARC 0.5, ARC 0.6}
  \begin{tabular}{lp{13cm}}
    User input:&\verb#(stderr=<string>)#\\
    GM input:&\verb#-"-#\\
    Example:&\verb#(stderr="myjob.err")#\\
  \end{tabular}

  The standard error file.

  \begin{tabular}{llp{10cm}}
    \hspace*{1cm}&\texttt{string} & file name, local to the computing element and
    relative to the session directory.\\
  \end{tabular}

  The standard error file should be listed as an \texttt{outputFiles}
  attribute; otherwise it will be forced to that list by the UI. If
  the standard error is not defined, UI assigns a name. If \texttt{join} is 
  specified with value "yes", UI adds \texttt{stderr} to the 
  pre-processed xRSL script with the same value as \texttt{stdout}.

  \hspace*{0.5cm}\phantomsection
  \index{XRSL:attribute:join}\addcontentsline{toc}{subsection}{join}
  \begin{shaded}
    \xrsl{join}
  \end{shaded}
  \versions{ARC 0.3, ARC 0.4, ARC 0.5, ARC 0.6}
  \begin{tabular}{lp{13cm}}
    User input:&\verb#(join="yes"|"no")#\\
    GM input:none\\
    Example:&\verb#(join="yes")#\\
  \end{tabular}

  If "yes", joins \texttt{stderr} and  \texttt{stdout} files into
  the \texttt{stdout} one. Default is \verb#no#.

  \hspace*{0.5cm}\phantomsection
  \index{XRSL:attribute:gmlog}\addcontentsline{toc}{subsection}{gmlog}
  \begin{shaded}
    \xrsl{gmlog}
  \end{shaded}
  \versions{ARC 0.3, ARC 0.4, ARC 0.5, ARC 0.6}
  \begin{tabular}{lp{13cm}}
    User input:&\verb#(gmlog=<string>)#\\
    GM input:&\verb#-"-#\\
    Example:&\verb#(gmlog="myjob.log")#\\
  \end{tabular}

  A name of the directory containing grid-specific diagnostics per job.

  \begin{tabular}{llp{10cm}}
    \hspace*{1cm}&\texttt{string} & a directory, local to the computing element and
    relative to the session directory\\
  \end{tabular}

%  The job log file should be listed as an \texttt{outputFiles} attribute;
%  otherwise it will be forced to that list by the UI.
	This directory is kept in the session directory to be
	available for retrieval (UI forces it to the list if
	\texttt{outputFiles})

  \hspace*{0.5cm}\phantomsection
  \index{XRSL:attribute:jobName}\addcontentsline{toc}{subsection}{jobName}
  \begin{shaded}
    \xrsl{jobName}
  \end{shaded}
  \versions{ARC 0.3, ARC 0.4, ARC 0.5, ARC 0.6}
  \begin{tabular}{lp{13cm}}
    User input:&\verb#(jobName=<string>)#\\
    GM input:&\verb#-"-#\\
    Example:&\verb#(jobName="My Job nr. 1")#\\
  \end{tabular}

  User-specified job name.

  \begin{tabular}{llp{10cm}}
    \hspace*{1cm}&\texttt{string} & job name\\
  \end{tabular}

  This name is meant for convenience of the user. It can be used to
  select the job while using the UI. It is also available through the
  Information System.

  \hspace*{0.5cm}\phantomsection
  \index{XRSL:attribute:ftpThreads}\addcontentsline{toc}{subsection}{ftpThreads}
  \begin{shaded}
    \xrsl{ftpThreads}
  \end{shaded}
  \versions{ARC 0.3, ARC 0.4, ARC 0.5, ARC 0.6}
  \begin{tabular}{lp{13cm}}
    User input:&\verb#(ftpThreads=<integer>)#\\
    GM input:&\verb#-"-#\\
    Example:&\verb#(ftpThreads="4")#\\
  \end{tabular}

  Defines how many parallel streams will be used by the GM during
  \texttt{gsiftp} and \texttt{http(s|g)}transfers of files.

  \begin{tabular}{llp{10cm}}
    \hspace*{1cm}&\texttt{integer} & a number from 1 to 10\\
  \end{tabular}

  If not specified, parallelism is not used.

  \hspace*{0.5cm}\phantomsection
  \index{XRSL:attribute:acl}\addcontentsline{toc}{subsection}{acl}
  \begin{shaded}
    \xrsl{acl}
  \end{shaded}
  \versions{ARC $\geq$ 0.5.12, ARC 0.6}

  \textit{Available in ARC v0.5.12 and higher}.

  \begin{tabular}{lp{13cm}}
    User input:&\verb#(acl=<xml>)#\\
    GM input:&\verb#-"-#\\
    Example:&\verb#(acl="<?xml version=""1.0""?>#\\
    &\verb#    <gacl version=""0.0.1""><entry><any-user></any-user>#\\
    &\verb#    <allow><write/><read/><list/><admin/></allow></entry></gacl>")#\\
  \end{tabular}



  Makes use of GACL~\cite{gacl} rules to list users who are allowed to
  access and control job in addition to job's owner. Access and
  control levels are specified per user. \texttt{any-user} tag refers
  to any user authorized at the execution cluster. To get more
  information about GACL please refer to \url{http://www.gridsite.org}.

  \begin{tabular}{llp{10cm}}
    \hspace*{1cm}&\texttt{xml} & a GACL-compliant XML string defining
    access control list\\
  \end{tabular}

  Following job control levels can be specified via \texttt{acl}:

  \begin{tabular}{llll}
    \hspace*{1cm}&\texttt{write}&--&\parbox[t]{11cm}{allows to modify contents of job data (job directory) and control job flow (cancel, clean, etc.)}\\
    \hspace*{1cm}&\texttt{read}&--&\parbox[t]{11cm}{allows to read content of job data
    (contents of job directory)}\\
    \hspace*{1cm}&\texttt{list}&--&\parbox[t]{11cm}{allows to list files available for
    the job (contents of job directory)}\\
    \hspace*{1cm}&\texttt{admin}&--&\parbox[t]{11cm}{allows to do everything --
    full equivalence to job ownership}\\
  \end{tabular}

  \hspace*{0.5cm}\phantomsection
  \index{XRSL:attribute:cluster}\addcontentsline{toc}{subsection}{cluster}
  \begin{shaded}
    \xrsl{cluster}
  \end{shaded}
  \versions{ARC 0.3, ARC 0.4, ARC 0.5, ARC 0.6}
  \begin{tabular}{lp{13cm}}
    User input:&\verb#(cluster=<string>)#\\
    GM input:&\verb#-"-#\\
    Example:&\verb#(cluster="cluster1.site2.org")#\\
  \end{tabular}

  The name of the execution cluster.

  \begin{tabular}{llp{10cm}}
    \hspace*{1cm}&\texttt{string} & known cluster name, or a substring of it\\
  \end{tabular}

  Use this attribute to explicitly force job submission to a cluster,
  or to avoid such. The job will not be submitted if the cluster does
  not satisfy other requirements of the job. Disjunct-requests of the
  kind \verb#(|(cluster="clus1")(cluster="clus2"))# are supported. To
  exclude a cluster, use \verb#(cluster!="clus3")#.

  \hspace*{0.5cm}\phantomsection
  \index{XRSL:attribute:queue}\addcontentsline{toc}{subsection}{queue}
  \begin{shaded}
    \xrsl{queue}
  \end{shaded}
  \versions{ARC 0.3, ARC 0.4, ARC 0.5, ARC 0.6}
  \begin{tabular}{lp{13cm}}
    User input:&\verb#(queue=<string>)#\\
    GM input:&\verb#-"-#\\
    Example:&\verb#(queue="pclong")#\\
  \end{tabular}
  
  The name of the remote batch queue. Use only if you are sure that
  the queue by this name does exist.

  \begin{tabular}{llp{10cm}}
    \hspace*{1cm}&\texttt{string} & known queue name\\
  \end{tabular}

  Use this attribute to explicitly force job submission to a queue.

  \hspace*{0.5cm}\phantomsection
  \index{XRSL:attribute:startTime}\addcontentsline{toc}{subsection}{startTime}
  \begin{shaded}
    \xrsl{startTime}
  \end{shaded}
  \versions{ARC 0.3, ARC 0.4, ARC 0.5, ARC 0.6}
  \begin{tabular}{lp{13cm}}
    User input:&\verb#(startTime=<time>)#\\
    GM input:&\verb#(startTime=<tttt>)#\\
    Example:&\verb#(startTime="2002-05-25 21:30")#\\
  \end{tabular}

  Time to start job processing by the Grid Manager, such as e.g. start downloading input files.

  \begin{tabular}{llp{10cm}}
    \hspace*{1cm}&\texttt{time} & time string, YYYY-MM-DD hh:mm:ss\\ 
    \hspace*{1cm}&\texttt{tttt} & time string, YYYYMMDDhhmmss[Z] (converted by the UI from
    \texttt{time})\\
  \end{tabular}
  
  \begin{framed}
   Actual job processing on a worker node starts depending on local scheduling mechanisms, but not sooner than \texttt{startTime}. 
   \end{framed}

  \hspace*{0.5cm}\phantomsection
  \index{XRSL:attribute:lifeTime}\addcontentsline{toc}{subsection}{lifeTime}
  \begin{shaded}
    \xrsl{lifeTime}
  \end{shaded}
  \versions{ARC 0.3, ARC 0.4, ARC 0.5, ARC 0.6}
  \begin{tabular}{lp{13cm}}
    User input:&\verb#(lifeTime=<time>)#\\
    GM input:&\verb#(lifeTime=<tttt>)#\\
    Example:&\verb#(lifeTime="60")#\\
  \end{tabular}

  Maximal time to keep job files (the session directory) on the
  gatekeeper upon job completion.

  \begin{tabular}{llp{10cm}}
    \hspace*{1cm}&\texttt{time} & time (days)\\
    \hspace*{1cm}&\texttt{tttt} & time (seconds, converted by the UI from \texttt{time})\\
  \end{tabular}

  Typical life time is 1 day (24 hours). Specified life time can not
  exceed local settings.

  \hspace*{0.5cm}\phantomsection
  \index{XRSL:attribute:notify}\addcontentsline{toc}{subsection}{notify}
  \begin{shaded}
    \xrsl{notify}
  \end{shaded}
  \versions{ARC 0.3, ARC 0.4, ARC 0.5, ARC 0.6}
  \begin{tabular}{lp{13cm}}
    User input:&\verb#(notify=<string> [string] ... )#\\
    GM input:&\verb#-"-#\\
    Example:&\verb#(notify="be your.name@your.domain.com")#\\
  \end{tabular}

  Request e-mail notifications on job status change.

  \begin{tabular}{llp{10cm}}
    \hspace*{1cm}&\texttt{string} & string of the format:
    \verb#[b][q][f][e][c][d] user1@domain1 [user2@domain2] ...#\\
    \hspace*{1cm}&&here flags indicating the job status are:\\
    \hspace*{1cm}&&\texttt{b} -- begin (PREPARING)\\
    \hspace*{1cm}&&\texttt{q} -- queued (INLRMS)\\
    \hspace*{1cm}&&\texttt{f} -- finalizing (FINISHING)\\
    \hspace*{1cm}&&\texttt{e} -- end (FINISHED)\\
    \hspace*{1cm}&&\texttt{c} -- cancellation (CANCELLED)\\
    \hspace*{1cm}&&\texttt{d} -- deleted (DELETED)\\
  \end{tabular}

  When no notification flags are specified, default value of ``\verb#eb#''
  will be used, i.e., notifications will be sent at the job's beginning
  and at its end.

  No more than 3 e-mail addresses per status change accepted.

  \hspace*{0.5cm}\phantomsection
  \index{XRSL:attribute:replicaCollection}\addcontentsline{toc}{subsection}{replicaCollection}
  \begin{shaded}
    \xrsl{replicaCollection}
  \end{shaded}
  \versions{ARC 0.3, ARC 0.4, ARC 0.5, ARC 0.6}
  \begin{tabular}{lp{13cm}}
    User input:&\verb#(replicaCollection=<URL>)#\\
    GM input:&\verb#-"-#\\
    Example:&\verb#(replicaCollection="ldap://grid.uio.no:389/lc=TestCollection,#\\
    &\verb#          rc=NorduGrid,nordugrid,dc=org")#\\
  \end{tabular}

  Location of a logical collection in the Replica Catalog.

  \begin{tabular}{llp{10cm}}
    \hspace*{1cm}&\texttt{URL} & LDAP directory specified as an URL (\verb#ldap://host[:port]/dn#)\\
  \end{tabular}

  \hspace*{0.5cm}\phantomsection
  \index{XRSL:attribute:rerun}\addcontentsline{toc}{subsection}{rerun}
  \begin{shaded}
    \xrsl{rerun}
  \end{shaded}
  \versions{ARC 0.3, ARC 0.4, ARC 0.5, ARC 0.6}
  \begin{tabular}{lp{13cm}}
    User input:&\verb#(rerun=<integer>)#\\
    GM input:&\verb#-"-#\\
    Example:&\verb#(rerun="2")#\\
  \end{tabular}

  Number of reruns (if a system failure occurs).

  \begin{tabular}{llp{10cm}}
    \hspace*{1cm}&\texttt{integer}  & an integer number\\
  \end{tabular}
    
  If not specified, the default is 0. Default maximal allowed value is
  5. Before 0.5.29 jobs were rerun automatically and only in case of
  failure in INLRMS state. Since 0.5.29 job may be rerun after failure
  in any state for which reruning has sense. To initiate rerun user
  has to use \texttt{ngresume} command.

  \hspace*{0.5cm}\phantomsection
  \index{XRSL:attribute:architecture}\addcontentsline{toc}{subsection}{architecture}
  \begin{shaded}
    \xrsl{architecture}
  \end{shaded}
  \versions{ARC 0.3, ARC 0.4, ARC 0.5, ARC 0.6}
  \begin{tabular}{lp{13cm}}
    User input:&\verb#(architecture=<string>)#\\
    GM input:&\\
    Example:&\verb#(architecture="i686")#\\
  \end{tabular}

  Request a specific architecture.

  \begin{tabular}{llp{10cm}}
    \hspace*{1cm}&\texttt{string} & architecture  (e.g., as produced by uname -a)\\
  \end{tabular}

  \hspace*{0.5cm}\phantomsection
  \index{XRSL:attribute:nodeAccess}\addcontentsline{toc}{subsection}{nodeAccess}
  \begin{shaded}
    \xrsl{nodeAccess}
  \end{shaded}
  \versions{ARC $\geq$ 0.3.24, ARC 0.4, ARC 0.5, ARC 0.6}
  \begin{tabular}{lp{13cm}}
    User input:&\verb#(nodeAccess="inbound"|"outbound")#\\
    GM input:&\\
    Example:&\verb#(nodeAccess="inbound")#\\
  \end{tabular}

  Request cluster nodes with inbound or outbound IP connectivity. If
  both are needed, a conjunct request should be specified.

  \hspace*{0.5cm}\phantomsection
  \index{XRSL:attribute:dryRun}\addcontentsline{toc}{subsection}{dryRun}
  \begin{shaded}
    \xrsl{dryRun}
  \end{shaded}
  \versions{ARC 0.3, ARC 0.4, ARC 0.5, ARC 0.6}
  \begin{tabular}{lp{13cm}}
    User input:&\verb#(dryRun="yes"|"no")#\\
    GM input:&\verb#-"-#\\
    Example:&\verb#(dryRun="yes")#\\
  \end{tabular}

  If "yes", do dry-run: RSL is parsed, but no job submission to LRMS
  is made. Should be used for xRSL and communication validation.

  \hspace*{0.5cm}\phantomsection
  \index{XRSL:attribute:rsl\_substitution}\addcontentsline{toc}{subsection}{rsl\_substitution}
  \begin{shaded}
    \xrsl{rsl\_substitution}
  \end{shaded}
  \versions{ARC 0.3, ARC 0.4, ARC 0.5, ARC 0.6}
  \begin{tabular}{lp{13cm}}
    User input:&\verb#(rsl_substitution=(<string1> <string2>))#\\
    GM input:&\verb#-"-#\\
    Example:&\verb#(rsl_substitution=("ATLAS" "/opt/atlas"))#\\
  \end{tabular}

  Substitutes \verb#<string2># with \verb#<string1># for
  \textbf{internal} RSL use.

  \begin{tabular}{llp{10cm}}
    \hspace*{1cm}&\texttt{string1} & new internal RSL variable\\
    \hspace*{1cm}&\texttt{string2} & any string, e.g., existing combination of
    variables or a path\\
  \end{tabular}
  
  Use this attribute to define variables that simplify xRSL editing,
  e.g. when same path is used in several values, typically in
  \texttt{inputFiles}. Only one pair per substitution is allowed. To
  request several substitution, concatenate such requests. Bear in
  mind that substitution must be defined \textbf{prior} to actual use
  of a new variable \texttt{string1}.

  After the substitution is defined, it should be used in a way
  similar to shell variables in scripts: enclosed in round brackets,
  preceded with a dollar sign, \textbf{without quotes}:\\
  \verb#(inputfiles=("myfile" $(ATLAS)/data/somefile))#

  Unlike the \texttt{environment} attribute, \texttt{rsl\_substitution}
  definition is only used by the client and is valid inside xRSL
  script. It can not be used to define environment or shell variable
  at the execution site.

  \hspace*{0.5cm}\phantomsection
  \index{XRSL:attribute:environment}\addcontentsline{toc}{subsection}{environment}
  \begin{shaded}
    \xrsl{environment}
  \end{shaded}
  \versions{ARC 0.3, ARC 0.4, ARC 0.5, ARC 0.6}
  \begin{tabular}{lp{13cm}}
    User input:&\verb#(environment=(<VAR> <string>) [(<VAR> <string>)] ... )#\\
    GM input:&\verb#-"-#\\
    Example:&\verb#(environment=("ATLSRC" "/opt/atlas/src")#\\
    &\verb#         ("ALISRC" "/opt/alice/src"))#\\
  \end{tabular}

  Defines execution shell environment variables.

  \begin{tabular}{llp{10cm}}
    \hspace*{1cm}&\texttt{VAR} & new variable name\\
    \hspace*{1cm}&\texttt{string} & any string, e.g., existing combination of
    variables or a path\\
  \end{tabular}

  Use this to define variables at an execution site. Unlike the
  \texttt{rsl\_substitution} attribute, it can not be used to define
  variables on the client side.

%%   \hspace*{0.5cm}\phantomsection
%%   \index{XRSL:attribute:jobtype}\addcontentsline{toc}{subsection}{jobtype}
%%   \begin{shaded}
%%     \xrsl{jobtype}
%%   \end{shaded}
%%   \versions{ARC $\geq$ 0.3.12, ARC 0.4}
%%   \begin{tabular}{lp{13cm}}
%%     User input:&\verb#(jobtype= "single"|"multiple"|"mpi"|"condor" )#\\
%%     Example:&\verb#(jobtype="mpi")#\\
%%   \end{tabular}

%%   Specifies a type of job (default is \texttt{single}).

  \hspace*{0.5cm}\phantomsection
  \addcontentsline{toc}{subsection}{count}\index{XRSL:attribute:count}
  \begin{shaded}
    \xrsl{count}
  \end{shaded}
  \versions{ARC $\geq$ 0.3.12, ARC 0.4, ARC 0.5, ARC 0.6}
  \begin{tabular}{lp{13cm}}
    User input:&\verb#(count=<integer>)#\\
    GM input:&\verb#-"-#\\
    Example:&\verb#(count="4")#\\
  \end{tabular}

  Specifies amount of sub-jobs to be submitted for parallel tasks.

  \hspace*{0.5cm}\phantomsection
  \addcontentsline{toc}{subsection}{jobreport}\index{XRSL:attribute:jobreport}
  \begin{shaded}
    \xrsl{jobreport}
  \end{shaded}
  \versions{ARC $\geq$ 0.3.34, ARC 0.4, ARC 0.5, ARC 0.6}
  \begin{tabular}{lp{13cm}}
    User input:&\verb#(jobreport=<URL>)#\\
    GM input:&\verb#-"-#\\
    Example:&\verb#(jobreport="https://grid.uio.no:8001/logger")#\\
  \end{tabular}

  Specifies an URL for a logging service to send reports about job
  to. The default is set up in the cluster configuration.

  \begin{tabular}{llp{10cm}}
    \hspace*{1cm}&\texttt{URL} & URL\\
  \end{tabular}

  It is up to a user to make sure the requested logging service
  accepts reports from the set of clusters she intends to use.

  \hspace*{0.5cm}\phantomsection
  \addcontentsline{toc}{subsection}{credentialserver}\index{XRSL:attribute:credentialserver}
  \begin{shaded}
    \xrsl{credentialserver}
  \end{shaded}
  \versions{ARC $\geq$ 0.5.45, ARC 0.6}
  \begin{tabular}{lp{13cm}}
    User input:&\verb#(credentialserver=<URL>)#\\
    GM input:&\verb#-"-#\\
    Example:&\verb#(credentialserver="myproxy://myproxy.nordugrid.org;username=user")#\\
  \end{tabular}

  Specifies an URL which Grid Manager may contact to renew/extend delegated
  proxy of job. Only MyProxy servers are supported.

  \begin{tabular}{llp{10cm}}
    \hspace*{1cm}&\texttt{URL} & URL of MyProxy server\\
  \end{tabular}

  It is up to a user to make sure the specified MyProxy server will accept
  requests from Grid Manager to renew expired credentials.
  \texttt{URL} may contain options \texttt{username} and \texttt{credname}
  to specify user name and credentials name which Grid Manager should pass
  to MyProxy server. If \texttt{username} is not specified DN of user 
  credentials is used instead.
  


  \section{GM-side attributes}
  \label{sec:gm}

  The following attributes are a part of the internal ARC job representation, and must be provided by ARC client tools and passed to
  the GM. Users may even specify them manually (\textbf{not
  advised!}), while developers of new ARC client tools and utilities must make sure these attributes are added to the user job description before it is submitted to a GM.

    \phantomsection
  \index{XRSL:attribute:sstdin}\addcontentsline{toc}{subsection}{sstdin}
  \hspace*{0.5cm}
  \begin{shaded}
    \xrsl{sstdin}
  \end{shaded}
  \versions{ARC 0.3, ARC 0.4, ARC $\leq$ 0.5.25}
  \begin{tabular}{lp{13cm}}
    User input:&\verb##\\
    GM input:&\verb#(sstdin=<filename>)#\\
    Example:&\verb#(sstdin="myinput.dat")#\\
  \end{tabular}

  Internal attribute for the standard input. Can also be spelled \texttt{stdinput}.
  Only needed for GRAM compatibility, not used by ARC as such.

  \begin{tabular}{llp{10cm}}
    \hspace*{1cm}&\texttt{filename} & standard input file name\\
  \end{tabular}

  \phantomsection
  \index{XRSL:attribute:action}\addcontentsline{toc}{subsection}{action}
  \hspace*{0.5cm}
  \begin{shaded}
    \xrsl{action}
  \end{shaded}
  \versions{ARC 0.3, ARC 0.4, ARC 0.5, ARC 0.6}
  \begin{tabular}{lp{13cm}}
    User input:&\verb##\\
    GM input:&\verb#(action="request"|"cancel"|"clean")#\\
    Example:&\verb#(action="request")#\\
  \end{tabular}

  Action to be taken by the gatekeeper: submit the job, cancel job
  execution, or clear the results of the job (also cancels the job).

  \phantomsection
  \index{XRSL:attribute:savestate}\addcontentsline{toc}{subsection}{savestate}
  \hspace*{0.5cm}
  \begin{shaded}
    \xrsl{savestate}
  \end{shaded}
  \versions{ARC 0.3, ARC 0.4, ARC $\leq$ 0.5.25}
  \begin{tabular}{lp{13cm}}
    User input:&\verb##\\
    GM input:&\verb#(savestate="yes"|"no")#\\
    Example:&\verb#(savestate="yes")#\\
  \end{tabular}

  If "yes", input RSL is stored in a temporary file at the
  gatekeeper. Must be always set as "yes" in the current implementation.
  Only needed for GRAM compatibility, not used by ARC as such.

  \phantomsection
  \index{XRSL:attribute:lrmstype}\addcontentsline{toc}{subsection}{lrmstype}
  \hspace*{0.5cm}
  \begin{shaded}
    \xrsl{lrmstype}
  \end{shaded}
  \versions{ARC $\leq$ 0.3.15}
  \begin{tabular}{lp{13cm}}
    User input:&\verb##\\
    GM input:&\verb#(lrmstype=<string>)#\\
    Example:&\verb#(lrmstype="pbs")#\\
  \end{tabular}

  LRMS type, indicating which submission script is to be invoked. 

  \begin{tabular}{llp{10cm}}
    \hspace*{1cm}&\texttt{string} & LRMS type \\
  \end{tabular}

  \phantomsection
  \index{XRSL:attribute:hostName}\addcontentsline{toc}{subsection}{hostName}
  \hspace*{0.5cm}
  \begin{shaded}
    \xrsl{hostName}
  \end{shaded}
  \versions{ARC 0.3, ARC 0.4, ARC 0.5, ARC 0.6}
  \begin{tabular}{lp{13cm}}
    User input:&\verb##\\
    GM input:&\verb#(hostname=<string>)#\\
    Example:&\verb#(hostName="grid.quark.lu.se")#\\
  \end{tabular}

  Submission host name. 

  \begin{tabular}{llp{10cm}}
    \hspace*{1cm}&\texttt{string} & host name, as passed by the UI\\
  \end{tabular}

  \phantomsection
  \index{XRSL:attribute:jobid}\addcontentsline{toc}{subsection}{jobid}
  \hspace*{0.5cm}
  \begin{shaded}
    \xrsl{jobid}
  \end{shaded}
  \versions{ARC 0.3, ARC 0.4, ARC 0.5, ARC 0.6}
  \begin{tabular}{lp{13cm}}
    User input:&\verb##\\
    GM input:&\verb#(jobid=<string>)#\\
    Example:&\verb#(jobid="grid.quark.lu.se:2119/jobmanager-ng/157111017133827")#\\
  \end{tabular}

  Unique job identification string, needed for cancellation and clean-up. 

  \begin{tabular}{llp{10cm}}
    \hspace*{1cm}&\texttt{string} & global job ID\\
  \end{tabular}

  It can also be provided during submission of the job and should be
  unique to a computing element (cluster).

  \phantomsection
  \index{XRSL:attribute:clientxrsl}\addcontentsline{toc}{subsection}{clientxrsl}
  \hspace*{0.5cm}
  \begin{shaded}
    \xrsl{clientxrsl}
  \end{shaded}
  \versions{ARC $\geq$ 0.3.28, ARC 0.4, ARC 0.5, ARC 0.6}
  \begin{tabular}{lp{13cm}}
    User input:&\verb##\\
    GM input:&\verb#(clientxrsl=<string>)#\\
    Example:&\verb#(clientxrsl="&(executable=/bin/echo)(arguments=boo)")#\\
  \end{tabular}

  Job description XRSL string as submitted by the user, before being pre-processed by the client. 

  \begin{tabular}{llp{10cm}}
    \hspace*{1cm}&\texttt{string} & original XRSL description submitted by the user\\
  \end{tabular}

  This attribute is added by the User Interface during pre-processing, and is used for job re-submission in order to repeat brokering and matchmaking.

  \phantomsection
  \index{XRSL:attribute:clientsoftware}\addcontentsline{toc}{subsection}{clientsoftware}
  \hspace*{0.5cm}
  \begin{shaded}
    \xrsl{clientsoftware}
  \end{shaded}
  \versions{ARC $\geq$ 0.3.35, ARC 0.4, ARC 0.5, ARC 0.6}
  \begin{tabular}{lp{13cm}}
    User input:&\verb##\\
    GM input:&\verb#(clientsoftware=<string>)#\\
    Example:&\verb#(clientsoftware="nordugrid-arc-0.5.39")#\\
  \end{tabular}

  Version of ARC client used to submit the job. 

  \begin{tabular}{llp{10cm}}
    \hspace*{1cm}&\texttt{string} & \\
  \end{tabular}

  This attribute is added by the User Interface during pre-processing.


  \section{Unsupported Globus RSL attributes}
  \label{sec:globus}

  The following \globus\ attributes are not supported by the ARC
  middleware. Whenever they are specified, eiher an error or a warning should be issued by the
  client tool, and the corresponding attribute should be ignored\footnote{Standard ARC client will throw an error, unless option ''-U'' is specified}.

  \subsection{Unsupported RSL 1.0 attributes}

  \begin{itemize}
    \item \verb#(resourceManagerContact=<string>)#
    \item \verb#(directory=<string>)#
    \item \verb#(maxCpuTime=<time>)#
    \item \verb#(maxWallTime=<time>)#
    \item \verb#(maxTime=<time>)#
    \item \verb#(maxMemory=<memory>)#
    \item \verb#(minMemory=<memory>)#
    \item \verb#(gramMyJob=independent|collective)#
    \item \verb#(project=<string>)#
    \item \verb#(hostCount=<number>)#
    \item \verb#(label=<string>)#
    \item \verb#(subjobCommsType=blocking-join|independent)#
    \item \verb#(subjobStartType=strict-barrier|loose-barrier|no-barrier)#
  \end{itemize}

  \subsection{Unsupported GRAM RSL attributes}

  \begin{itemize}
    \item \verb#(directory=<string>)#
    \item \verb#(fileCleanUp=<array>)#
    \item \verb#(fileStageIn=<array>)#
    \item \verb#(fileStageInShared=<array>)#
    \item \verb#(fileStageOut=<array>)#
    \item \verb#(gassCache=<path>)#
    \item \verb#(gramMyJob=independent|collective)#
    \item \verb#(hostCount=<number>)#
    \item \verb#(jobType=<string>)#
    \item \verb#(libraryPath=<path>)#
    \item \verb#(maxCpuTime=<time>)#
    \item \verb#(maxWallTime=<time>)#
    \item \verb#(maxTime=<time>)#
    \item \verb#(maxMemory=<memory>)#
    \item \verb#(minMemory=<memory>)#
    \item \verb#(project=<string>)#
    \item \verb#(remoteIoUrl=<string>)#
    \item \verb#(scratchDir=<string>)#
  \end{itemize}

  \chapter{XRSL and JSDL}
  \label{sec:jsdl}

    	
  ARC has adopted JSDL~\cite{jsdl} as second format for describing job
  requests\footnote{Available in releases $\geq$ 0.5.37, ARC 0.6, ARC 0.8}. ARC also adds own extensions to JSDL. Below supported JSDL 
  elements are listed together with corresponding XRSL
  attributes. Here \textit{jsdl} stands for jsdl namespace,
  \textit{jsdlPOSIX} -- for POSIX extensions, and \textit{jsdlARC} --
  for ARC extensions.

\begin{longtable}{|l|l|p{5cm}|}
\hline 
\textbf{XRSL attribute} & \textbf{JSDL element} & \textbf{Comments} \\ \hline 
executable & jsdlPOSIX:Executable & \\ \hline 
cpuTime & \parbox[t]{5.5cm}{
    jsdlPOSIX:CPUTimeLimit \\
    jsdl:Resources \\
    \hspace*{0.5cm}jsdl:TotalCpuTime\\
    \hspace*{0.5cm}jsdl:IndividualCpuTime} & \\ \hline 
arguments & jsdlPOSIX:Argument & \\ \hline 
inputFiles &  \parbox[t]{5.5cm}{
    jsdl:DataStaging \\
    \hspace*{0.5cm}jsdl:Source \\
    \hspace*{1cm}jsdl:FileName \\
    \hspace*{1cm}jsdl:URI} &
    jsdl:DataStaging without URI in Source are treated like user-uploadable files \\ \hline
outputFiles & \parbox[t]{5.5cm}{
    jsdl:DataStaging \\
    \hspace*{0.5cm}jsdl:Target \\
    \hspace*{1cm}jsdl:FileName \\
    \hspace*{1cm}jsdl:URI} &
    jsdl:DataStaging without Target and Source or without URI in Target are treated like user-downloadable files \\ \hline 
wallTime & jsdlPOSIX:WallTimeLimit & \\ \hline 
memory &  \parbox[t]{5.5cm}{
    jsdlPOSIX:MemoryLimit \\
    jsdl:Resources \\
    \hspace*{0.5cm}jsdl:TotalPhysicalMemory \\
    \hspace*{0.5cm}jsdl:IndividualPhysicalMemory} & \\ \hline
disk & \parbox[t]{5.5cm}{
    jsdl:Resources \\
    \hspace*{0.5cm}jsdl:TotalDiskSpace \\
    \hspace*{0.5cm}jsdl:IndividualDiskSpace} & \\ \hline
stdin & jsdlPOSIX:Input & \\ \hline
stdout & jsdlPosix:Output & \\ \hline
stderr & jsdlPOSIX:Error & \\ \hline
jobName & jsdl:JobIdentification & \\ \hline
architecture & \parbox[t]{5.5cm}{
    jsdl:Resources \\
    \hspace*{0.5cm}jsdl:CPUArchitectureName \\
    \hspace*{1cm}jsdl:CPUArchitecture} & \\ \hline
count &  \parbox[t]{5.5cm}{
    jsdl:Resources \\
    \hspace*{0.5cm}jsdl:TotalCpuCount \\
    \hspace*{0.5cm}jsdl:IndividualCpuCount} & \\ \hline
executables &  \parbox[t]{5.5cm}{
    jsdl:DataStaging \\
    \hspace*{0.5cm} jsdl-arc:IsExecutable} & \\ \hline
gridTime &  \parbox[t]{5.5cm}{
    jsdl:Resources \\
    \hspace*{0.5cm} jsdl-arc:GridTimeLimit} & \\ \hline
runTimeEnvironment & \parbox[t]{5.5cm}{
    jsdl:Resources \\
    \hspace*{0.5cm}jsdl-arc:RunTimeEnvironment} &
    Only "Exact" condition with one Version is supported so far \\ \hline
middleware &  \parbox[t]{5.5cm}{
    jsdl:Resources \\
    \hspace*{0.5cm}jsdl-arc:Middleware} &
    Only "Exact" condition with one Version is supported so far \\ \hline
gmlog &  \parbox[t]{5.5cm}{
    jsdl-arc:LocalLogging \\
    \hspace*{0.5cm}jsdl-arc:Directory} & \\ \hline
acl &  \parbox[t]{5.5cm}{
    jsdl-arc:AccessControl \\
    \hspace*{0.5cm}Type=GACL \\
    \hspace*{0.5cm}Content} & \\ \hline
cluster &  \parbox[t]{5.5cm}{
    jsdl:Resources \\
    \hspace*{0.5cm}jsdl-arc:CandidateTarget \\
    \hspace*{1cm}jsdl-arc:HostName} & \\ \hline
queue &  \parbox[t]{5.5cm}{
    jsdl:Resource \\
    \hspace*{0.5cm}jsdl-arc:CandidateTarget \\
    \hspace*{1cm}jsdl-arc:QueueName} & \\ \hline
startTime & jsdl-arc:ProcessingStartTime & \\ \hline
lifeTime & \parbox[t]{5.5cm}{
    jsdl:Resources \\
    \hspace*{0.5cm}jsdl-arc:SessionFileTime} & \\ \hline
notify & \parbox[t]{5.5cm}{
    jsdl-arc:Notify \\
    \hspace*{0.5cm}jsdl-arc:Type=Email} & \\ \hline
rerun & jsdl-arc:Reruns & \\ \hline
jobreport & \parbox[t]{5.5cm}{
    jsdl-arc:RemoteLogging \\
    \hspace*{0.5cm}jsdl-arc:URL} & \\ \hline
credentialserver & \parbox[t]{5.5cm}{
    jsdl-arc:CredentialServer \\
    \hspace*{0.5cm}jsdl-arc:URL} & \\ \hline
\end{longtable}

Conditional requests are supported through recursive jsdl:JobDescription - 
jsdl:JobDescriptin inside jsdl:JobDescription.


  \appendix

  \chapter{Examples}
  \label{sec:examples}

  \section{User-side xRSL script}

  \begin{verbatim}
&
(* test run: if "yes", only submits RSL without actual job start *)
     (dryRun="no")
(* some local variables defined for further convenience *)
     (rsl_substitution=("TOPDIR" "/home/johndoe"))
     (rsl_substitution=("NGTEST" $(TOPDIR)/ngtest))
     (rsl_substitution=("BIGFILE" "/scratch/johndoe/100mb.tmp"))
(* some environment variables, to be used by the job *)
     (environment=("ATLAS" "/opt/atlas") ("CERN" "/cern"))
(* the main executable file to be staged in and submitted to the PBS *)
     (executable="checkall.sh")
(* the arguments for the executable above *)
     (arguments="pal")
(* files to be staged in before the execution *)
     (inputFiles = ("be_kaons" "")
     ("file1" gsiftp://grid.uio.no$(TOPDIR)/remfile.txt)
      ("bigfile.dat" $(BIGFILE) )    )
(* files to be given executable permissions after staging in *)
     (executables="be_kaons")
(* files to be staged out after the execution *)
     (outputFiles=
        ("file1" "gsiftp://grid.tsl.uu.se/tmp/file1.tmp")
        ("100mb.tmp" "rls://rls.nordugrid.org:39281/test/bigfile")
        ("be_kaons.hbook"  gsiftp://ce1.grid.org$(NGTEST)/kaons.hbook)   )
(* user-specified job name *)
     (jobName="NGtest")
(* standard input file *)
     (stdin="myinput.dat")
(* standard output file *)
     (stdout="myoutput.dat")
(* standard error file *)
     (stderr="myerror.dat")
(* GM logs directory name *)
     (gmlog="gmlog")
(* flag whether to merge stdout and stderr *)
     (join="no")
(* request e-mail notification on status change *)
     (notify="bqfe john.doe@gmail.com jane.doe@mail.org")
(* maximal CPU time required for the job, minutes for PBS*)
     (CpuTime="60")
(* maximal time for the session directory to exist on the remote node, days *)
     (lifeTime="7")
 (* memory required for the job, per rank, Mbytes *)
     (Memory="200")
(* wall time to start job processing *)
     (startTime="2002-04-28 17:15:00")
(* disk space required for the job, Mbytes *)
     (Disk="500")
(* required architecture of the execution node *)
     (architecture="i686")
(* required run-time environment *)
     (runTimeEnvironment="APPS/HEP/Atlas-1.1")
(* number of re-runs, in case of a system failure *)
     (rerun="2")
  \end{verbatim}

  \section{GM-side xRSL script}

Note that the client tool must expand all the \texttt{rsl\_substitution}
values and add double quotes to \textbf{all} strings. Client tool is also expected to insert \texttt{queue} attribute in case such is missing and to add checksum values for input files to be uploaded by the client. Standard ARC client also
streamlines blank spaces. 

Comment lines are removed from the GM-side xRSL; below they are shown
to explain details.

  \begin{verbatim}
&
(* saves RSL in a temporary file if "yes" *)
     ("savestate" = "yes" )
(* job submission to be performed if action is "request" *)
     ("action" = "request" )
(* submission host name *)
     ("hostname" = "ce1.grid.org" )
(* client software version *)
     ("clientsoftware" = "nordugrid-arc-0.6.0.3" )
(* walltime value added by the client, in seconds *)
     ("walltime" = "3600" )
(* test run: if "yes", only submits RSL without actual job start *)
     ("dryRun" = "no" )
(* some local variables defined for further convenience *)
     ("rsl_substitution" = ("TOPDIR" "/home/johndoe" ) )
     ("rsl_substitution" = ("NGTEST" "/home/johndoe/ngtest" ) )
     ("rsl_substitution" = ("BIGFILE" "/scratch/johndoe/100mb.tmp" ) )
(* some environment variables, to be used by the job *)
     ("environment" = ("ATLAS" "/opt/atlas" ) ("CERN" "/cern" ) )
(* executable *)
     ("executable" = "checkall.sh" )
(* arguments *)
     ("arguments" = "pal" )
(* files to be staged in before the execution *)
     ("inputfiles" = ("checkall.sh" "279320" )
        ("myinput.dat" "39806" )
        ("be_kaons" "8807" )
        ("file1" "gsiftp://grid.uio.no/home/johndoe/remfile.txt" )
        ("bigfile.dat" "104857600" )
     )
(* files to be given executable permissions after staging in *)
     ("executables" = "checkall.sh" "be_kaons" )
(* files to be staged out after the execution *)
     ("outputfiles" = ("file1" "gsiftp://grid.tsl.uu.se/tmp/file1.tmp" )
        ("100mb.tmp" "rls://rls.nordugrid.org:39281/test/bigfile" )
        ("be_kaons.hbook" "gsiftp://ce1.grid.org/home/johndoe/ngtest/kaons.hbook" )
        ("myoutput.dat" "" )
        ("myerror.dat" "" )
    )
(* user-specified job name *)
     ("jobName" = "NGtest" )
(* standard input file *)
     ("stdin" = "myinput.dat" )
(* standard output file *)
     ("stdout"="myoutput.dat" )
(* standard error file *)
     ("stderr" = "myerror.dat" )
(* flag whether to merge stdout and stderr *)
     ("join" = "no" )
(* request e-mail notification on status change *)
     ("notify" = "bqfe john.doe@gmail.com jane.doe@mail.org" )
(* specific queue to submit the job *)
     ("queue" = "atlas" )
(* CPU time required for the job, converted into seconds *)
     ("cputime" = "3600" )
(* maximal time for the session directory to exist on the remote node, seconds *)
     ("lifetime" = "604800" )
(* memory required for the job, per rank, Mbytes *)
     ("memory" = "200" )
(* wall time to start job processing *)
     ("startTime" = "20020128171500" )
(* disk space required for the job, Mbytes *)
     ("disk" = "500" )
(* required architecture of the execution node *)
     ("architecture" = "i686" )
(* required run-time environment *)
     ("runtimeenvironment" = "APPS/HEP/Atlas-1.1" )
(* number of re-runs, in case of a system failure *)
     ("rerun" = "2" )
(* original client xRSL with expanded string substitutions; shortened here *)
     ("clientxrsl" = "&(""dryrun"" = ""no"" )(... )(""rerun"" = ""2"" )" )
  \end{verbatim}

  \section{JSDL script}

  \begin{verbatim}

<?xml version="1.0" encoding="UTF-8"?>
<jsdl:JobDefinition
  xmlns="http://schemas.ggf.org/jsdl/2005/11/jsdl"
  xmlns:posix="http://schemas.ggf.org/jsdl/2005/11/jsdl-posix"
  xmlns:arc="http://www.nordugrid.org/ws/schemas/jsdl-arc">
  <JobDescription>
    <!-- Generic job description (jobname in XRSL) -->
    <JobIdentification>
        <JobName>JSDL Hello World job</JobName>
        <Description>
          This is a simple Hello World job test
        </Description>
    </JobIdentification>
    <!-- Standard application (same as executable, arguments in XRSL) -->
    <Application>
        <ApplicationName>Echo</ApplicationName>
        <posix:POSIXApplication>
          <posix:Executable>/bin/echo</posix:Executable>
          <posix:Argument>Hello world</posix:Argument>
        </posix:POSIXApplication>
    </Application>
    <!-- Get the GM diagnostics (same as gmlog in XRSL) -->
    <arc:LocalLogging>
	      <arc:Directory>gmlog</arc:Directory>
    </arc:LocalLogging>
 </JobDescription>
</JobDefinition>


  \end{verbatim}


  \chapter{JSDL-ARC schema}
  \label{sec:jsdl-arc-schema}

  
  \begin{verbatim}

<?xml version="1.0" encoding="UTF-8"?>
<xsd:schema xmlns:xsd="http://www.w3.org/2001/XMLSchema"
	    xmlns="http://www.nordugrid.org/ws/schemas/jsdl-arc"
	    xmlns:jsdl-arc="http://www.nordugrid.org/ws/schemas/jsdl-arc"
	    targetNamespace="http://www.nordugrid.org/ws/schemas/jsdl-arc">

 <xsd:simpleType name="GMState_Type">
  <xsd:restriction base="xsd:string">
   <xsd:enumeration value="ACCEPTED"/>
   <xsd:enumeration value="PREPARING"/>
   <xsd:enumeration value="SUBMIT"/>
   <xsd:enumeration value="INLRMS"/>
   <xsd:enumeration value="FINISHING"/>
   <xsd:enumeration value="FINISHED"/>
   <xsd:enumeration value="DELETED"/>
   <xsd:enumeration value="CANCELING"/>
  </xsd:restriction>
 </xsd:simpleType>

 <xsd:complexType name="Version_Type">
  <xsd:sequence>
   <xsd:element name="UpperExclusive" type="xsd:string" minOccurs="0"/>
   <xsd:element name="LowerExclusive" type="xsd:string" minOccurs="0"/>
   <xsd:element name="Exact" type="xsd:string" minOccurs="0" maxOccurs="unbounded"/>
   <xsd:element name="Exclusive" type="xsd:boolean" minOccurs="0"/>
  </xsd:sequence>
 </xsd:complexType>

 <xsd:simpleType name="SessionType_Type">
  <xsd:documentation> For jsdl:Resources_Type </xsd:documentation>
  <!-- xsd:element ref="SessionType" minOccurs="0"/ -->
  <xsd:restriction base="xsd:string">
   <xsd:enumeration value="INTERNAL"/>
   <xsd:enumeration value="LIMITED"/>
   <xsd:enumeration value="READONLY"/>
   <xsd:enumeration value="FULL"/>
  </xsd:restriction>
 </xsd:simpleType>

 <xsd:simpleType name="IsExecutable_Type">
  <xsd:documentation> For jsdl:DataStaging_Type (default: false) </xsd:documentation>
  <!-- xsd:element ref="IsExecutable" minOccurs="0"/ -->
  <xsd:restriction base="xsd:boolean"/>
 </xsd:simpleType>

 <xsd:simpleType name="FileParameters_Type">
  <xsd:documentation> For jsdl:DataStaging_Type </xsd:documentation>
  <!-- xsd:element ref="IsExecutable" minOccurs="0"/ -->
  <xsd:restriction base="xsd:string"/>
 </xsd:simpleType>

 <xsd:simpleType name="JoinOutputs_Type">
  <xsd:documentation> For jsdl:JobDescription_Type (default: false) </xsd:documentation>
  <!-- xsd:element ref="JoinOutputs" minOccurs="0"/ -->
  <xsd:restriction base="xsd:boolean"/>
 </xsd:simpleType>

 <xsd:simpleType name="Reruns_Type">
  <xsd:documentation> For jsdl:JobDescription_Type (default: false) </xsd:documentation>
  <!-- xsd:element ref="Reruns" minOccurs="0"// -->
  <xsd:restriction base="xsd:integer"/>
 </xsd:simpleType>

 <xsd:complexType name="RunTimeEnvironment_Type">
  <xsd:documentation> For jsdl:Resources_Type </xsd:documentation>
  <!-- xsd:element ref="RunTimeEnvironment" minOccurs="0" maxOccurs="unbounded"/ -->
  <xsd:sequence>
   <xsd:element name="Name" type="xsd:string"/>
   <xsd:element name="Version" type="Version_Type" minOccurs="0"/>
  </xsd:sequence>
 </xsd:complexType>

 <xsd:complexType name="Middleware_Type">
  <xsd:documentation> For jsdl:Resources_Type </xsd:documentation>
  <!-- xsd:element ref="Middleware" minOccurs="0" maxOccurs="unbounded"/ -->
  <xsd:sequence>
   <xsd:element name="Name" type="xsd:string"/>
   <xsd:element name="Version" type="Version_Type" minOccurs="0"/>
  </xsd:sequence>
 </xsd:complexType>

 <xsd:complexType name="RemoteLogging_Type">
  <xsd:documentation> For jsdl:JobDescription_Type </xsd:documentation>
  <!-- xsd:element ref="RemoteLogging" minOccurs="0" maxOccurs="3"/ -->
  <xsd:sequence>
   <xsd:element name="URL" minOccurs="1" maxOccurs="1" type="xsd:anyURI"/>
  </xsd:sequence>
 </xsd:complexType>

 <xsd:complexType name="CredentialServer_Type">
  <xsd:documentation> For jsdl:JobDescription_Type </xsd:documentation>
  <!-- xsd:element ref="CredentialServer" minOccurs="0"/ -->
  <xsd:sequence>
   <xsd:element name="URL" minOccurs="1" maxOccurs="1" type="xsd:anyURI"/>
  </xsd:sequence>
 </xsd:complexType>

 <xsd:complexType name="LocalLogging_Type">
  <xsd:documentation> For jsdl:JobDescription_Type </xsd:documentation>
  <!-- xsd:element ref="LocalLogging" minOccurs="0" maxOccurs="1"/ -->
  <xsd:sequence>
   <xsd:element name="Directory" minOccurs="1" maxOccurs="1" type="xsd:string"/>
  </xsd:sequence>
 </xsd:complexType>

 <xsd:simpleType name="AccessControlType_Type">
  <xsd:restriction base="xsd:string">
   <xsd:enumeration value="GACL"/>
  </xsd:restriction>
 </xsd:simpleType>
 <xsd:complexType name="AccessControl_Type">
  <xsd:documentation> For jsdl:JobDescription_Type </xsd:documentation>
  <!-- xsd:element ref="AccessControl" minOccurs="0"/ -->
  <xsd:sequence>
   <xsd:element name="OwnerAlwaysAllowed" type="xsd:boolean" minOccurs="0"/>
   <xsd:element name="Type" type="AccessControlType_Type" minOccurs="0"/>
   <xsd:element name="Content" minOccurs="0" type="xsd:string"/>
  </xsd:sequence>
 </xsd:complexType>

 <xsd:simpleType name="NotificationType_Type">
  <xsd:restriction base="xsd:string">
   <xsd:enumeration value="Email"/>
  </xsd:restriction>
 </xsd:simpleType>
 <xsd:complexType name="Notify_Type">
  <xsd:documentation> For jsdl:JobDescription_Type </xsd:documentation>
  <!-- xsd:element ref="Notify" minOccurs="0" maxOccurs="3"/ -->
  <xsd:sequence>
   <xsd:element name="Type" type="NotificationType_Type" minOccurs="0"/>
   <xsd:element name="Endpoint" minOccurs="0" type="xsd:string"/>
   <xsd:element name="State" minOccurs="1" maxOccurs="unbounded" type="GMState_Type"/>
  </xsd:sequence>
 </xsd:complexType>

 <xsd:simpleType name="SessionLifeTime_Type">
  <xsd:documentation> For jsdl:Resources_Type </xsd:documentation>
  <!-- xsd:element ref="SessionLifeTime" minOccurs="0" maxOccurs="1"/ -->
  <xsd:restriction base="xsd:long"/>
 </xsd:simpleType>

 <xsd:simpleType name="GridTimeLimit_Type">
  <xsd:documentation> For jsdl:Resources_Type </xsd:documentation>
  <!-- xsd:element ref="GridTimeLimit" minOccurs="0" maxOccurs="1"/ -->
  <xsd:restriction base="xsd:positiveInteger"/>
 </xsd:simpleType>
 
 <xsd:complexType name="CandidateTarget_Type">
  <xsd:documentation> For jsdl:Resources_Type </xsd:documentation>
  <!-- xsd:element ref="jsdl-arc:CandidateTarget" minOccurs="0" maxOccurs="1"/ -->
  <xsd:sequence>
   <xsd:element name="HostName" minOccurs="0" type="xsd:string"/>
   <xsd:element name="QueueName" minOccurs="0" type="xsd:string"/>
  </xsd:sequence>
 </xsd:complexType>

 <xsd:simpleType name="Time_Type">
  <xsd:documentation> For jsdl:JobDescription_Type </xsd:documentation>
  <!-- xsd:element ref="ProcessingStartTime" minOccurs="0" maxOccurs="1"/ -->
  <xsd:restriction base="xsd:dateTime"/>
 </xsd:simpleType>

 <!--=========================================================-->
 <xsd:element name="IsExecutable" type="IsExecutable_Type"/>
 <xsd:element name="FileParameters" type="FileParameters_Type"/>
 <xsd:element name="RunTimeEnvironment" type="RunTimeEnvironment_Type"/>
 <xsd:element name="Middleware" type="Middleware_Type"/>
 <xsd:element name="RemoteLogging" type="RemoteLogging_Type"/>
 <xsd:element name="LocalLogging" type="LocalLogging_Type"/>
 <xsd:element name="AccessControl" type="AccessControl_Type"/>
 <xsd:element name="Notify" type="Notify_Type"/>
 <xsd:element name="SessionLifeTime" type="SessionLifeTime_Type"/>
 <xsd:element name="SessionType" type="SessionType_Type"/>
 <xsd:element name="JoinOutputs" type="JoinOutputs_Type"/>
 <xsd:element name="Reruns" type="Reruns_Type"/>
 <xsd:element name="CredentialServer" type="CredentialServer_Type"/>
 <xsd:element name="GridTimeLimit" type="GridTimeLimit_Type"/>
 <xsd:element name="CandidateTarget" type="CandidateTarget_Type"/>
 <xsd:element name="ProcessingStartTime" type="Time_Type"/>

</xsd:schema>

  \end{verbatim}



  \bibliography{grid}
  \printindex

\end{document}
